\section{Optimal proxies for nearly balanced odd populations}\label{sec:odd_asym}
The strict singleton comparison of Theorem \ref{odd:thm:odd} has a small gap for nearly balanced populations. This section determines that gap through fourth order and the location of the globally maximizing tilt through third order, for every fixed sample size and fixed success imbalance. The cases $m=1$ and $m=2$ agree with the Bernoulli formula \eqref{eq:bernproxy} and the three-point specialization in Section \ref{sec:twodraw}. The theorem gives the expansion for arbitrary fixed $m$, together with uniqueness of the global optimizer and analytic remainders.

Let $N$ be odd, let $s\ge1$ be a fixed odd integer, and put
\begin{equation}\label{eq:oddparams}
 k=\frac{N-s}{2},\qquad H_N\sim\Hyp(N,k,m),\qquad
 R_{N,s,m}=\frac{4N}{m(N-m)}\sigma^2_{\mathrm{opt}}(H_N).
\end{equation}
Here $m\ge1$ is fixed while $N\to\infty$ runs through odd integers larger than $s+2m$. The prescribed centering gives a Bernoulli-success fraction $p=(1-s/N)/2$. The regime includes the most nearly balanced success count when $s=1$.

\begin{theorem}[Fourth-order sharp-proxy expansion]\label{thm:odd_expansion}
Fix $m\ge1$ and a positive odd integer $s$. For all sufficiently large odd $N$, the function
\[
 t\longmapsto\frac{2}{t^2}\log\E\exp\{t(H_N-mk/N)\},
\]
extended continuously at $t=0$, has a unique global maximizer $t_{N,s,m}>0$. The maximizer and the optimal normalized proxy have the expansions
\begin{align}
 t_{N,s,m}&=\frac{4s}{N}+\frac{4s(m-1)}{N^2}
 +\frac{4s(6m+s^2-6)}{3N^3}+O_{m,s}(N^{-4}),\label{eq:tiltexp}\\
 R_{N,s,m}
 &=1+\frac1N+\frac{1-s^2/3}{N^2}
 +\frac{1-(2m-1)s^2/3}{N^3}\notag\\
 &\quad+\frac{1-(4m^2-2m-1)s^2/3-4s^4/45}{N^4}
 +O_{m,s}(N^{-5}).\label{eq:Rexp}
\end{align}
The remainder is analytic in $N^{-1}$ in a neighborhood of zero for every fixed $m,s$.
\end{theorem}

\begin{corollary}[Fourth-order singleton dominance]\label{cor:fixedm_gap}
For every fixed integer $m\ge2$, let $H_N\sim\Hyp(N,(N-1)/2,m)$, $N$ odd, and let
\[
 L_N=\frac{1}{(N-1)\operatorname{artanh}(1/N)}.
\]
Then
\begin{equation}\label{eq:fixedgap}
 \boxed{\quad L_N-R_{N,1,m}
 =\frac{2(m-1)}{3N^3}
 +\frac{2(m-1)(2m+1)}{3N^4}+O_m(N^{-5}).\quad}
\end{equation}
In particular, $R_{N,1,m}<L_N$ for all sufficiently large odd $N$, for every fixed $m\ge2$. More generally, if $s\ge3$ is fixed and odd, then
\begin{equation}\label{eq:offs_gap}
 L_N-R_{N,s,m}=\frac{s^2-1}{3N^2}+O_{m,s}(N^{-3})>0
\end{equation}
for all sufficiently large odd $N$.
\end{corollary}

\subsection{Cumulants from factorial moments}\label{sec:cumulants_odd}
We compute the cumulants entering the globally optimized expansion in Theorem \ref{thm:odd_expansion}. Write $\varepsilon=N^{-1}$. For $r\le m$, the factorial moments of $H_N$ satisfy
\begin{equation}\label{eq:factorial_moments}
 \E(H_N)_{\underline r}
 =\frac{(m)_{\underline r}(k)_{\underline r}}
 {(N)_{\underline r}},\qquad k=\frac{N-s}{2}.
\end{equation}
The ordinary moments up to order six are obtained by the Stirling transformation
\[
 \E H_N^q=\sum_{r=0}^q\left\{\begin{matrix}q\\r\end{matrix}\right\}
 \E(H_N)_{\underline r}.
\]
The expressions remain valid when $r>m$, with zero falling factorials. Let $\kappa_j(\varepsilon)$ denote the $j$th cumulant of $H_N$. Elementary simplification of the first six moments yields
\begin{align}
 \kappa_2(e)&=\frac m4\left[1-(m-1)e-(m+s^2-1)e^2
 +(m-1)(s^2-1)(e^3+e^4)\right]+O_{m,s}(e^5),\label{eq:kappa2exp}\\
 \kappa_3(e)&=\frac{ms}{4}\left[e-3(m-1)e^2
 +(2m^2-9m-s^2+7)e^3\right]+O_{m,s}(e^4),\label{eq:kappa3exp}\\
 \kappa_4(e)&=-\frac m8+\frac{m(m-1)}2e
 +\frac{m(-6m^2+19m+4s^2-13)}8e^2+O_{m,s}(e^3),\label{eq:kappa4exp}\\
 \kappa_5(e)&=-\frac{ms}{2}e+O_{m,s}(e^2),\qquad
 \kappa_6(e)=\frac m4+O_{m,s}(e).\label{eq:kappa56exp}
\end{align}
Here $e=\varepsilon$. All expansions follow from factorial moments through order six and the cumulant recursion
\[
 \kappa_r=\mu_r-\sum_{j=1}^{r-1}\binom{r-1}{j-1}\kappa_j\mu_{r-j},
 \qquad \mu_r=\E H^r.
\]
For the variance and third cumulant one can verify the first two identities directly from
\begin{align}
 \kappa_2(\varepsilon)
 &=\frac m4(1-s^2\varepsilon^2)\frac{1-m\varepsilon}{1-\varepsilon},\label{eq:exactk2}\\
 \kappa_3(\varepsilon)
 &=\kappa_2(\varepsilon)\,s\varepsilon
 \frac{1-2m\varepsilon}{1-2\varepsilon}.
 \label{eq:exactk3}
\end{align}
Formula \eqref{eq:kappa4exp} follows from \eqref{eq:factorial_moments} for $r=1,2,3,4$, using
\[
 \kappa_4=\E(H-\E H)^4-3\kappa_2^2.
\]
All these identities are rational in $\varepsilon$ for fixed $m,s$ and are valid in a neighborhood of the origin. The zero-parameter limit is $H_0\sim\mathrm{Binomial}(m,1/2)$, for which odd cumulants vanish and $\kappa_4(0)=-m/8$.

\subsection{Localization and uniqueness of the maximizing tilt}
Let
\[
 Q_\varepsilon(t)=\begin{cases}
  2\psi_\varepsilon(t)/t^2,&t\ne0,\\
  \kappa_2(\varepsilon),&t=0.
 \end{cases}
\]
For every fixed $m,s$, the finite probability generating polynomial depends analytically on $(t,\varepsilon)$ near $(0,0)$, and its value at $(0,0)$ is one. Consequently $Q_\varepsilon$ is jointly real analytic in a neighborhood of $(0,0)$. At $\varepsilon=0$,
\begin{equation}\label{eq:q0}
 Q_0(t)=\frac{2m}{t^2}\log\cosh(t/2),\qquad Q_0(0)=m/4.
\end{equation}
The strict inequality $\log\cosh x<x^2/2$ for $x\ne0$ implies
\begin{equation}\label{eq:qstrict}
 Q_0(t)<Q_0(0)\qquad(t\ne0).
\end{equation}
Furthermore, $|H_N-\E H_N|\le m$ implies $0\le Q_\varepsilon(t)\le2m/|t|$ for $t\ne0$, uniformly in the admissible populations. On every compact interval, $Q_\varepsilon$ converges uniformly to $Q_0$, including at zero by analytic continuation. Combining this compact convergence with the uniform tail bound shows that every global maximizer of $Q_\varepsilon$ tends to zero as $\varepsilon\downarrow0$.

Near $(0,0)$, the cumulant expansion gives
\begin{equation}\label{eq:Qexpand}
 Q_\varepsilon(t)=\kappa_2(\varepsilon)
 +\frac{\kappa_3(\varepsilon)}3t
 +\frac{\kappa_4(\varepsilon)}{12}t^2
 +\frac{\kappa_5(\varepsilon)}{60}t^3
 +\frac{\kappa_6(\varepsilon)}{360}t^4+O_{m,s}(t^5).
\end{equation}
At the limiting law, $\kappa_5(0)=0$, and
\[
 \partial_t^2Q_0(0)=\frac{\kappa_4(0)}6=-\frac m{48}<0.
\]
Thus $Q_\varepsilon$ is strictly concave on some fixed neighborhood of zero for all sufficiently small $\varepsilon$. The derivative at zero is $\kappa_3(\varepsilon)/3>0$ for $\varepsilon>0$ small. By the analytic implicit-function theorem, there exists a unique critical point $t(\varepsilon)$ in that neighborhood; it is analytic in $\varepsilon$, and the localization argument establishes that it is the unique global maximizer. Differentiation of the critical point equation at zero gives
\[
 t'(0)=-\frac{\partial_\varepsilon\partial_t Q(0,0)}{\partial_t^2Q(0,0)}
 =-\frac{ms/12}{-m/48}=4s.
\]
This proves the leading term of \eqref{eq:tiltexp} and the positivity of the maximizing tilt. The next subsection computes the remaining coefficients.

\subsection{Evaluation through fourth order}
The analytic critical-point equation can be expanded coefficient by coefficient:
\[
 0=\frac{\kappa_3(e)}3+\frac{\kappa_4(e)}6t
 +\frac{\kappa_5(e)}{20}t^2+\frac{\kappa_6(e)}{90}t^3+O_{m,s}(t^4).
\]
Set $t=a e+b e^2+c e^3+O(e^4)$. The coefficients of $e,e^2,e^3$, in order, give
\[
 a=4s,\qquad b=4s(m-1),\qquad c=\frac{4s(6m+s^2-6)}3.
\]
The terms omitted from the derivative contribute $O_{m,s}(e^4)$ because the cumulants are analytic and $t=O(e)$. This establishes the full tilt expansion in \eqref{eq:tiltexp}.

For the optimal value through order $e^4$, it suffices to use
\[
 \widehat t=4se+4s(m-1)e^2.
\]
Indeed $t_* -\widehat t=O(e^3)$ and $Q_e'(t_*)=0$, so Taylor's theorem around the true optimizer gives $Q_e(t_*)-Q_e(\widehat t)=O(e^6)$. With $\kappa_5(0)=0$ and $\kappa_6(0)=m/4$, substitution into \eqref{eq:Qexpand} yields
\begin{align*}
 \sigma^2_{\mathrm{opt}}(H_N)
 &=\kappa_2+\frac{\kappa_3}3\widehat t
 +\frac{\kappa_4}{12}\widehat t^2
 +\frac{\kappa_5}{60}\widehat t^3
 +\frac{\kappa_6}{360}\widehat t^4+O_{m,s}(e^5)\\
 &=\frac m4\Bigl[1-(m-1)e+(1-m-s^2/3)e^2\\
 &\hspace{14mm}-(m-1)(1+s^2/3)e^3\\
 &\hspace{14mm}+\{1-m-(2m^2-m-1)s^2/3-4s^4/45\}e^4\Bigr]
 +O_{m,s}(e^5).
\end{align*}
Multiplying by $4/[m(1-me)]$ proves the fourth-order proxy expansion \eqref{eq:Rexp}. The analytic global optimizer supplies the Taylor remainder for each fixed $m,s$.

The exact singleton benchmark is
\begin{equation}\label{eq:singletonL}
 L_N=\frac{1}{(N-1)\operatorname{artanh}(1/N)}
 =1+e+\frac23e^2+\frac23e^3+\frac{26}{45}e^4+O(e^5).
\end{equation}
Subtracting \eqref{eq:Rexp} from this expression proves Corollary \ref{cor:fixedm_gap}.

\subsection{Finite-size interpretation}
The $N^{-1}$ coefficient in \eqref{eq:Rexp} is one. The $N^{-2}$ coefficient records the success imbalance through $-s^2/3$ and is common to all fixed sample sizes. For $s=1$, the first sample-size contribution is the order-$N^{-3}$ term $-2(m-1)/3$ relative to the singleton law. Thus \eqref{eq:fixedgap} quantifies the separation between the optimized proxies. The maximizing tilt has leading term $4s/N$.

\section{Uniform extremizers at fixed sample size}\label{sec:uniform_fixed}
At each fixed sample size, the following results maximize the proxy over every success count and compute its local separation from the Bernoulli envelope.

\subsection{A two-parameter analytic family}
Fix an integer $m\ge1$. Write $e=N^{-1}$, $d=1-2k/N$, $p=(1-d)/2$, and $q=1-p$. For $i=0,\ldots,m$, define
\begin{equation}\label{eq:analytic_masses}
 P_i(e,d)=\binom mi
 \frac{\prod_{r=0}^{i-1}(p-re)\prod_{r=0}^{m-i-1}(q-re)}
 {\prod_{r=0}^{m-1}(1-re)}.
\end{equation}
Empty products equal one. For physical parameters this is the hypergeometric probability mass function. Vandermonde's falling-factorial identity gives $\sum_iP_i=1$ and $\sum_i iP_i=mp$. In a rectangle around $(e,d)=(0,0)$, every factor is positive and the masses are analytic. At $e=0$ the law is $\mathrm{Binomial}(m,p)$. Define
\[
 \psi(e,d,t)=\log\sum_{i=0}^mP_i(e,d)e^{t(i-mp)},\qquad
 Q(e,d,t)=\frac{2\psi(e,d,t)}{t^2},
\]
with $Q(e,d,0)=\partial_t^2\psi(e,d,0)$. Division by $t^2$ is removable because $\psi(e,d,0)=\partial_t\psi(e,d,0)=0$. Let
\[
 \mathcal R_m(e,d)=\frac{4}{m(1-me)}\max_{t\in\mathbb R}Q(e,d,t).
\]

\begin{lemma}[Joint analytic optimizer]\label{lem:joint_optimizer}
For each fixed $m$, there is a rectangle around $(0,0)$ on which $Q(e,d,\cdot)$ has a unique global maximizer $\tau_m(e,d)$. Both $\tau_m$ and $\mathcal R_m$ are real analytic. Moreover,
\begin{align}
 \tau_m(e,-d)&=-\tau_m(e,d),&
 \mathcal R_m(e,-d)&=\mathcal R_m(e,d),\label{eq:joint_symmetry}\\
 \tau_m(e,0)&=0,&
 \mathcal R_m(e,0)&=\frac1{1-e}.\label{eq:balanced_exact}
\end{align}
\end{lemma}
\begin{proof}
At $(e,d)=(0,0)$, $Q(t)=2m\log\cosh(t/2)/t^2$ has its unique maximum at zero and $Q_{tt}(0)=-m/48$. For nearby parameters, $|H-mp|\le m$ gives $0\le Q(t)\le2m/|t|$. Choose $T$ so large that this tail bound lies strictly below $m/4$. On $[-T,T]$, joint continuity, including at $t=0$, yields uniform convergence to the limiting objective. Its strict gap outside any neighborhood of zero forces every maximizing point into that neighborhood. Joint analyticity and $Q_{tt}(0,0,0)<0$ then give a fixed smaller neighborhood where $Q_{tt}<0$. The analytic implicit-function theorem applied to $Q_t=0$ gives one critical point there, which the localization argument identifies as the unique global maximum.

The identity $P_i(e,-d)=P_{m-i}(e,d)$ implies $Q(e,-d,t)=Q(e,d,-t)$, proving \eqref{eq:joint_symmetry}. At $d=0$ the objective is even in $t$, so uniqueness gives $\tau_m(e,0)=0$. Its value is the exact variance
\[
 \partial_t^2\psi(e,0,0)=\frac m4\frac{1-me}{1-e},
\]
which proves \eqref{eq:balanced_exact}.
\end{proof}

\begin{theorem}[Uniform fixed-sample maximum]\label{thm:uniform_fixed}
For each fixed integer $m\ge1$, there exists $N_0(m)$ such that, for every odd $N\ge N_0(m)$,
\begin{equation}\label{eq:fixed_all_k}
 \max_{1\le k<N}\frac{4N}{m(N-m)}
 \sigma^2_{\mathrm{opt}}(H_{N,k,m})=R_{N,1,m}.
\end{equation}
The maximizing success counts are exactly $k=(N-1)/2$ and $k=(N+1)/2$. If $m\ge2$, this maximum is strictly below $L_N$. Consequently, for every fixed $M\ge2$, all sample sizes $2\le m\le M$ and all success counts are simultaneously below $L_N$ once $N$ is sufficiently large.
\end{theorem}
\begin{proof}
At $e=0$, the Bernoulli proxy formula gives
\[
 \mathcal R_m(0,d)=g(d):=\frac d{\operatorname{artanh}d},\qquad g(0)=1.
\]
In particular $g(d)=1-d^2/3+O(d^4)$. The even analytic function from Lemma \ref{lem:joint_optimizer} therefore has the form
\[
 \mathcal R_m(e,d)=\frac1{1-e}-d^2A_m(e,d^2),\qquad A_m(0,0)=\frac13,
\]
where $A_m$ is analytic. Choose $d_0>0$ and $e_0>0$ inside its analytic rectangle such that
\[
 A_m(e,d^2)+d^2\partial_2A_m(e,d^2)>0
 \quad(0\le e\le e_0,\ 0\le d\le d_0).
\]
It follows that $\partial_d\mathcal R_m(e,d)<0$ for $0<d\le d_0$. Thus, among odd-population success counts in this interval, the maximum is uniquely attained at $d=1/N=e$.

For the remaining success counts, the elementary-symmetric-mean inequality gives
\begin{equation}\label{eq:maclaurin_mgf}
 \E e^{tH_{N,k,m}}\le(q+pe^t)^m.
\end{equation}
Indeed, $\E e^{tH}$ is the normalized $m$th elementary symmetric polynomial of $N$ positive weights, $k$ equal to $e^t$ and the others equal to one. At fixed sum, averaging a pair of weights increases that polynomial by $(a-b)^2e_{m-2}(\mathrm{rest})/4\ge0$. Repeatedly averaging the largest and smallest weights decreases the sum of their squares by $(a-b)^2/2$; consequently their difference tends to zero, and the weights converge to their common mean. The elementary symmetric polynomial never decreases during this procedure. The case $m=1$ is equality. This proves \eqref{eq:maclaurin_mgf}. After centering, the Bernoulli proxy formula yields
\[
 \frac{4N}{m(N-m)}\sigma^2_{\mathrm{opt}}(H_{N,k,m})
 \le\frac{g(d)}{1-me}\le\frac{g(d_0)}{1-me}\quad(d\ge d_0).
\]
Since $g(d_0)<1$, this last bound is below one for sufficiently small $e>0$, whereas $\mathcal R_m(e,e)=1+e+O_m(e^2)>1$. Reflection covers negative $d$. This proves \eqref{eq:fixed_all_k} and the equality cases. The strict comparison with $L_N$ follows from Corollary \ref{cor:fixedm_gap}. Taking the largest of finitely many thresholds proves the assertion for $2\le m\le M$.
\end{proof}


\subsection{Exact local separation at finite population size}
The polynomial proof of Theorem \ref{thm:curvature} also applies to the normal form with integer $j\ge1$ and real $h,r\ge0$. In this extension, the coefficients in \eqref{eq:Q} use Gamma functions for factorial ratios. Their ratio, the differential equation, and the stationary identity remain unchanged; Jacobi orthogonality holds for real $h,r\ge0$, and the positivity certificate uses only $S=h+r\ge0$. The cases $j\ge3$ and the two boundary layers therefore prove the same curvature bound. This observation allows a continuous success-count perturbation at $k=N/2$, including odd $N$.

\begin{theorem}[Exact imbalance coefficient]\label{thm:exact_imbalance}
Fix integers $N\ge5$ and $1\le m\le\lfloor N/2\rfloor$. Extend the hypergeometric masses to real $k=N(1-d)/2$ by \eqref{eq:analytic_masses}, with $e=1/N$. For all sufficiently small $|d|$, the optimal normalized proxy has a unique maximizing tilt, analytic in $d$, and
\begin{align}
 \tau_m(1/N,d)&=T_{N,m}d+O_{N,m}(d^3),\label{eq:finite_tilt}\\
 \mathcal R_m(1/N,d)&=\frac N{N-1}-A_{N,m}d^2+O_{N,m}(d^4),\label{eq:finite_imbalance}
\end{align}
where, writing $D_{N,m}=N^2-N-3m(N-m)>0$,
\begin{align}
 T_{N,m}&=\frac{4(N-2m)(N-3)(N-1)}{(N-2)D_{N,m}},\label{eq:finite_T}\\
 A_{N,m}&=\frac N{3(N-1)}+
 \frac{2N^2(N-4)(m-1)(N-m-1)}
 {3(N-2)^2(N-1)D_{N,m}}.\label{eq:finite_A}
\end{align}
For $m\ge2$, $A_{N,m}>A_{N,1}$. Thus the Bernoulli envelope strictly dominates the $m$-draw proxy for every sufficiently small nonzero real imbalance.
\end{theorem}
\begin{proof}
At $d=0$, use the real-margin extension with $j=m$ and $h=r=N/2-m\ge0$. The curvature envelope equals the variance
\[
 B=\frac{m(N-m)}{4(N-1)}.
\]
Its twice-integrated bound is strict at every nonzero $t$, by the analytic-gap argument in Corollary \ref{cor:mgf}; hence $Q(1/N,0,t)$ has its unique global maximum at zero. The masses remain positive for $|d|$ sufficiently small. Their finite support gives uniform tail control, so the localization argument in Lemma \ref{lem:joint_optimizer} applies at this fixed $N$.

Factorial moments give, at $d=0$,
\begin{align*}
 \kappa_2(d)&=B(1-d^2),\\
 \kappa_3(d)&=B\frac{N-2m}{N-2}d+O(d^3),\\
 \kappa_4(0)&=-\frac{BD_{N,m}}{2(N-3)(N-1)}.
\end{align*}
Since $m(N-m)\le N^2/4$, $D_{N,m}\ge N(N-4)/4>0$. Consequently $Q_{tt}(1/N,0,0)=\kappa_4(0)/6<0$, and the analytic implicit-function theorem gives a unique analytic maximizing tilt. If $b=B(N-2m)/(N-2)$ and $c=\kappa_4(0)$, the stationary equation gives $T_{N,m}=-2b/c$, proving \eqref{eq:finite_T}. Substitution into
\[
 Q(d,t)=B(1-d^2)+\frac{bd}{3}t+\frac c{12}t^2
 +O_{N,m}(d^4)\qquad(t=O(d))
\]
gives
\[
 \max_tQ(d,t)=B-\left(B+\frac{b^2}{3c}\right)d^2+O_{N,m}(d^4).
\]
Reflection makes the optimized value even in $d$, so the stated remainder follows. Multiplying by $4N/[m(N-m)]$ yields
\[
 A_{N,m}=\frac N{N-1}\left[
 1-\frac{2(N-2m)^2(N-3)(N-1)}{3(N-2)^2D_{N,m}}\right].
\]
Subtracting $N/[3(N-1)]$ and putting over a common denominator gives exactly the positive second term of \eqref{eq:finite_A}. The singleton proxy is $\mathcal R_1(1/N,d)=N g(d)/(N-1)$, so its quadratic coefficient is $A_{N,1}=N/[3(N-1)]$. The asserted local domination follows from the strict coefficient difference.
\end{proof}

\subsection{Contraction on compact success-proportion intervals}
The first finite-population correction gives a comparison away from the balanced point as well.

\begin{theorem}[Uniform Bernoulli envelope at fixed sample size]\label{thm:compact_contraction}
Fix $m\ge2$ and $0<\eta<1/2$. There exists $N_1(m,\eta)$ such that, for all $N\ge N_1(m,\eta)$ and $\eta\le k/N\le1-\eta$,
\begin{equation}\label{eq:compact_envelope}
 \sigma^2_{\mathrm{opt}}(H_{N,k,m})
 \le\frac{m(N-m)}{N-1}\,
 \sigma^2_{\mathrm{opt}}(\Bern(k/N)).
\end{equation}
Equality holds when $k=N/2$; the inequality is strict otherwise.
\end{theorem}
\begin{proof}
By reflection let $d\in[0,1-2\eta]$. At $e=0$ the binomial proxy has its unique maximizing tilt $t_0(d)=4\operatorname{artanh}d$, with continuous value $t_0(0)=0$. Its second derivative there is negative: for $d>0$, the Bernoulli proof in Appendix \ref{app:bernoulli} gives $G''(t_0)=c-pq>0$ and $Q''(t_0)=-2mG''(t_0)/t_0^2<0$; at $d=0$ the derivative is $-m/48$. Compactness, the uniform tail bound, and the strict uniqueness at $e=0$ allow the localization and implicit-function arguments to be applied uniformly on this interval. Thus $\mathcal R_m(e,d)$ is analytic on a neighborhood of its compact parameter interval for sufficiently small $|e|$.

For $p,q$ in this compact interval, expanding \eqref{eq:analytic_masses} gives
\[
 P_i(e,d)=\binom mi p^iq^{m-i}\left[
 1+\frac e2\left\{m(m-1)-\frac{i(i-1)}p
 -\frac{(m-i)(m-i-1)}q\right\}\right]+O_{m,\eta}(e^2).
\]
If $Z=q+pe^t$, binomial factorial moments therefore yield
\begin{equation}\label{eq:first_order_psi}
 \psi(e,d,t)=m\log Z-mpt+
 \frac{em(m-1)}2\left[1-\frac{pe^{2t}+q}{Z^2}\right]+O_{m,\eta}(e^2)
\end{equation}
on compact $t$ intervals. Analytic divisibility by $t^2$ permits the same expansion for $Q$ at zero. Evaluation at the unique limiting maximizer gives, with $a=\log(q/p)=2\operatorname{artanh}d$,
\[
 \mathcal R_m(e,d)=g(d)+e\left[m g(d)-(m-1)\frac{d^2}{a^2pq}\right]+O_{m,\eta}(e^2).
\]
The continuous interpretation is used at $d=0$. Since $\mathcal R_1(e,d)=g(d)/(1-e)$, their difference vanishes identically at $e=0$ and also at $d=0$ for every small $e$. It is even in $d$. Joint analyticity consequently factors it exactly as
\[
 \mathcal R_1(e,d)-\mathcal R_m(e,d)=e d^2 B_m(e,d),
\]
where $B_m$ is continuous on the compact interval and
\[
 B_m(0,d)=\frac{(m-1)g(d)}{d^2}
 \left[\frac{d}{(1-d^2)\operatorname{artanh}d}-1\right],
 \qquad B_m(0,0)=\frac{2(m-1)}3.
\]
For $d>0$, integration gives $\operatorname{artanh}d<d/(1-d^2)$, so $B_m(0,d)>0$. Its compact minimum is positive. Uniform continuity then gives $B_m(e,d)>0$ for all sufficiently small $e>0$. This proves the normalized comparison, including strictness, and multiplying through gives \eqref{eq:compact_envelope}.
\end{proof}

\begin{remark}[Relation to the finite-population envelope]
Theorem \ref{thm:compact_contraction} gives a local analytic comparison for both parities on compact success-proportion intervals. For odd populations, Theorem \ref{thm:main_laplace} gives the strict proxy comparison at every finite sample size in the canonical triangular domain. Theorem \ref{thm:uniform_fixed} further identifies the maximizing success count at each fixed sample size for all sufficiently large odd $N$.
\end{remark}


\section{Pointwise curvature and its weighted integral}\label{sec:curvature_obstruction}
The exact odd-population constant defines the factor
\[
 \rho_N=\frac{N-1}{N}L_N=\frac{1}{N\operatorname{artanh}(1/N)}<1.
\]
The corresponding pointwise inequality takes the form
\begin{equation}\label{eq:falseoddcurve}
 \psi_{N,k,m}''(t)\le\rho_N\frac{m(N-m)}{4(N-1)}.
\end{equation}
The next proposition gives an exact tilted variance above the right side of \eqref{eq:falseoddcurve}.

\begin{proposition}[Exact exceedance of the constant pointwise bound]\label{prop:obstruction}
For $N=15$, $k=7$, $m=3$, and $t=\log(7/6)$,
\begin{equation}\label{eq:obs_exact}
 \frac{4(N-1)}{m(N-m)}\psi_{N,k,m}''(t)
 =\frac{307887776}{308178025}
 >\frac{675}{676}>\rho_{15}.
\end{equation}
This nearest-balanced population violates \eqref{eq:falseoddcurve} at the stated tilt.
\end{proposition}

\begin{proof}
With $z=7/6$, the tilted hypergeometric weights for $j=0,1,2,3$ are proportional to
\[
 w_j=\binom7j\binom8{3-j}z^j.
\]
They are rational, so both the mean and the variance are rational. Direct summation gives
\[
 \E_tH=\frac{26313}{17555},\qquad
 \operatorname{Var}_t(H)=\frac{197927856}{308178025}.
\]
Since $m(N-m)/[4(N-1)]=9/14$, the first equality in \eqref{eq:obs_exact} follows by multiplying the variance by $14/9$. The first strict inequality follows from
\[
 \frac{307887776}{308178025}-\frac{675}{676}
 =\frac{111969701}{208328344900}>0.
\]
Finally, the power series for $\operatorname{artanh}$ at $x=1/15$ yields
\[
 \operatorname{artanh}(1/15)>\frac1{15}+\frac1{3\cdot15^3},
 \quad\text{hence}\quad
 \rho_{15}<\left(1+\frac1{3\cdot15^2}\right)^{-1}=\frac{675}{676}.
\]
Combining these two inequalities proves \eqref{eq:obs_exact}.
\end{proof}

\begin{proposition}[Integrated curvature identity]\label{prop:integral}
For every centered random variable with finite moment generating function, writing $\psi(t)=\log\E e^{tX}$ gives
\begin{equation}\label{eq:curv_integral}
 \frac{2\psi(t)}{t^2}=2\int_0^1(1-u)\psi''(ut)\,du,
 \qquad t\ne0.
\end{equation}
The proxy objective is therefore a weighted average of curvature along the segment from zero to the chosen tilt.
\end{proposition}
\begin{proof}
Twice integrate $\psi''$ using $\psi(0)=\psi'(0)=0$ and substitute $s=ut$; the formula is valid for either sign of $t$.
\end{proof}

Theorem \ref{odd:thm:envelope} bounds the weighted integral \eqref{eq:curv_integral} through the mean-dependent envelope of Lemma \ref{odd:lem:barrier}. Proposition \ref{prop:obstruction} identifies a point of this variance curve above the constant threshold in \eqref{eq:falseoddcurve}.
